\section{Datos y Metodología}
\label{sec:datos-metodologia}

\subsection{Fuente y estructura de los datos}

El conjunto de datos utilizado (\texttt{Actividad\_2\_Data.csv}) se deriva del histórico
global de temperatura superficial compilado por \citet{berkeleyearth}. Cada registro
corresponde a la temperatura promedio mensual observada en una ciudad, con las siguientes
columnas: \texttt{dt} (fecha, primer día del mes), \texttt{Average\-Temperature} (temperatura
promedio en grados Celsius), \texttt{AverageTemperatureUncertainty} (incertidumbre de la
medición), \texttt{City}, \texttt{Country}, \texttt{Latitude} y \texttt{Longitude}. El
archivo cubre \DatasetCityCount{} ciudades entre 1743 y 2013.

\subsection{Limpieza de datos}

Del total de registros originales, se descartaron aquellos sin temperatura observada
(\texttt{Average\-Temperature} nulo), producto de vacíos en el proceso de medición histórica.
Tras esta depuración quedaron \DatasetRowCount{} registros válidos distribuidos entre las
\DatasetCityCount{} ciudades. Las fechas se tipificaron como \texttt{datetime} y, para cada
ciudad, los registros se ordenaron cronológicamente antes de cualquier análisis posterior.

\subsection{Ventana de análisis para Bogotá}

Al reindexar la serie de Bogotá a una frecuencia mensual regular se identificó un
\textbf{vacío continuo de 20 años} (enero de 1830 a junio de 1850, 246 meses consecutivos sin
registro), además de un puñado de meses aislados sin dato en otros periodos. Rellenar un
vacío de esa magnitud mediante interpolación habría implicado \emph{fabricar} dos décadas de
tendencia y estacionalidad inexistentes en los datos, en lugar de simplemente completar
huecos puntuales. Por esta razón, el análisis y el modelado se restringen al periodo
continuo y confiable de la serie: \textbf{\AnalysisStartYear{}--\AnalysisEndYear{}}, que
comprende \BogotaObsCount{} observaciones mensuales sin vacíos relevantes. Los pocos meses
aislados sin registro dentro de esta ventana sí se completan mediante interpolación lineal
en el tiempo, una operación razonable cuando el hueco es de uno o dos meses.

\subsection{Metodología}

\paragraph{Análisis gráfico.} Se grafica la serie histórica completa de Bogotá y se compara
su perfil climático promedio mes a mes contra ciudades representativas de climas distintos
(templado boreal, desértico/subtropical, tropical monzónico y templado del hemisferio sur),
independientemente del rango de años disponible en cada una. Adicionalmente, se construye un
boxplot mensual para inspeccionar la dispersión y detectar valores atípicos por mes.

\paragraph{Descomposición.} Se aplica una descomposición clásica de tipo aditivo
($Y_t = T_t + S_t + R_t$) con periodo estacional de 12 meses, mediante la función
\texttt{seasonal\_decompose} de la librería \texttt{statsmodels}
\citep{seabold2010statsmodels}. Se opta por el modelo aditivo en lugar del multiplicativo
porque la amplitud del ciclo estacional de temperatura no muestra una relación proporcional
clara con el nivel de la serie.

\paragraph{Modelado y pronóstico.} Se divide la serie en un conjunto de entrenamiento y uno
de prueba, reservando los últimos \TestSizeMonths{} meses para validación. Sobre el conjunto
de entrenamiento se aplica la prueba de Dickey-Fuller aumentada (ADF)
\citep{dickeyfuller1979} para evaluar estacionariedad, y se inspeccionan las funciones de
autocorrelación (ACF) y autocorrelación parcial (PACF) como insumo para acotar el orden del
modelo. El orden final $(p,d,q)\times(P,D,Q)_{12}$ de un modelo SARIMA se selecciona mediante
una búsqueda en grilla que minimiza el criterio de información de Akaike (AIC), ajustando los
modelos con \texttt{SARIMAX} de \texttt{statsmodels}. El modelo seleccionado se evalúa sobre
el conjunto de prueba mediante el error cuadrático medio (RMSE) y el error absoluto medio
(MAE); posteriormente se reajusta sobre la serie completa para proyectar
\FutureHorizonMonths{} meses hacia el futuro, junto con su intervalo de confianza del 95\,\%.

\subsection{Herramientas y reproducibilidad}

Todo el análisis se implementó en Python con tipado estático estricto (verificado con
\texttt{mypy --strict}), utilizando \texttt{pandas} \citep{mckinney2010pandas},
\texttt{statsmodels} \citep{seabold2010statsmodels}, \texttt{scikit-learn} y
\texttt{matplotlib}. Un único script (\texttt{src/run\_pipeline.py}) ejecuta de punta a punta
la carga y limpieza de datos, el análisis gráfico, la descomposición y el modelado,
regenerando de forma determinística tanto las figuras como los valores numéricos citados en
este informe.
